\chapter{Options Market Structure}\label{ch:m1:options-market-structure}

A retail customer taps ``buy 10 calls at the market''. The order never
reaches an order book in the ordinary sense. Her broker sends it to a
\omterm{def:m1:retail-flow-and-wholesaling:wholesaler}{wholesaler}; the \omterm{def:m1:retail-flow-and-wholesaling:wholesaler}{wholesaler} submits it to an exchange together with its own
order to take the other side at a guaranteed price; the exchange broadcasts a
request for responses and waits, for a period that may be as short as a
tenth of a second, while the largest market-making firms in the world decide
whether to offer a cent better. Then the order is divided: part to those who
improved, a reserved part to the firm that brought it, the rest shared by
size. Every step of this is written in an exchange rule. The options market
is where the mechanisms of the last five chapters (wholesaling, \omterm{def:m1:matching-algorithms-and-implied-spreads:prorata}{pro-rata
allocation}, market-maker privileges, fee models) are combined most densely,
and where a firm's position in the structure decides more of its profit than
its models do.

\section{A quote-driven market}

\begin{definition}[Options market maker and quoting obligation]\label{def:m1:options-market-structure:mm}
An \emph{options market maker}\index{options market maker} is a firm
registered with an exchange to quote options classes. It accepts a
\emph{quoting obligation}\index{quoting obligation}, typically to quote two
sides in a set percentage of the series of its classes, for a set percentage
of the trading day, within a maximum width, and receives privileges in
return: lower fees, allocation entitlements, the right to send quotes in
bulk.
\end{definition}

With two thousand series on one share (\cref{ch:m1:option-contracts-and-exchanges}),
most of which trade a few times a day or not at all, natural buyers and
sellers almost never meet. Nearly every displayed price is a \omterm{def:m1:what-a-trading-firm-does:market-maker}{market maker}'s
quote, computed from the share price and a volatility surface and refreshed
whenever either moves. The book is a projection of a model, and the economic
questions are who may quote, what they must quote, and how an incoming order
is divided among them.

\begin{definition}[Quote protection]\label{def:m1:options-market-structure:protection}
\emph{Quote protection}\index{quote protection} denotes the exchange-provided
mechanisms that limit a \omterm{def:m1:what-a-trading-firm-does:market-maker}{market maker}'s exposure across its many quotes: a
risk monitor that cancels all of a firm's quotes in a class when its fills
within a short window exceed a set number, bulk cancel messages, and
kill switches.
\end{definition}

A \omterm{def:m1:what-a-trading-firm-does:market-maker}{market maker} quoting two thousand series is, economically, offering one
bet on the share in two thousand forms. An informed trader who lifts a
hundred of them in a millisecond has traded a hundred times against the same
stale opinion. Without protection, quotes would be a fraction of their size.

\section{Allocation: customers first, then size}

\begin{definition}[Customer priority]\label{def:m1:options-market-structure:customer}
Under \emph{customer priority}\index{customer priority} resting orders of
\emph{priority customers} (broadly, the public, as opposed to \omterm{def:m1:the-sell-side:sell-side}{broker-dealers}
and professional traders) at the best price are filled before any
professional interest at that price, whatever their time of arrival. The
remainder is allocated among professionals, usually pro rata on size, often
after an entitlement for a designated \omterm{def:m1:what-a-trading-firm-does:market-maker}{market maker}.
\end{definition}

\begin{omfigure}
\begin{tikzpicture}
\begin{axis}[width=11.5cm,height=5.2cm,ybar,bar width=8pt,
  ylabel={lots received},
  symbolic x coords={MM1,C1,DMM,F,C2,MM2},xtick=data,
  xticklabels={{MM1: 100},{C1: 10},{DMM: 100},{F: 50},{C2: 5},{MM2: 250}},
  tick label style={font=\small},label style={font=\small},xticklabel style={font=\footnotesize\strut},
  ymin=0,ymax=66,ymajorgrids,grid style={gray!25},enlarge x limits=0.1,
  nodes near coords,nodes near coords style={font=\scriptsize},
  legend columns=3,legend style={font=\footnotesize,at={(0.5,-0.2)},anchor=north,draw=none,fill=none,/tikz/every even column/.append style={column sep=8pt}},
  table/col sep=comma]
\addplot[fill=omExo!50,draw=omExo] table[x=order,y=pro_rata]{figdata/markets-1/24-options-market-structure/rules.csv};
\addplot[fill=omDef!70,draw=omDef] table[x=order,y=customer]{figdata/markets-1/24-options-market-structure/rules.csv};
\addplot[fill=omThm!60,draw=omThm] table[x=order,y=entitled]{figdata/markets-1/24-options-market-structure/rules.csv};
\legend{plain pro rata,customer priority,customer priority and a 40\% entitlement}
\end{axis}
\end{tikzpicture}

\omcaption{An order for 115 lots meets a level of six resting orders, listed
in time priority with their sizes: two customers (C), three \omterm{def:m1:what-a-trading-firm-does:market-maker}{market makers},
one of them designated (DMM), and a firm (F). Customers are filled in full
under priority; the entitlement doubles the designated maker's share at the
expense of the other professionals. Data: the chapter's
build.}\label{fig:m1:options-market-structure:rules}
\end{omfigure}

The rule is a transfer. Customers, whose orders carry the least information,
are worth the most to trade against; giving them priority attracts them to
the exchange, and their flow is then the prize that the professionals'
allocation rules divide. Exchanges that want to attract \omterm{def:m1:what-a-trading-firm-does:market-maker}{market makers}'
quotes offer entitlements and pro rata; exchanges that want to attract
takers offer \omterm{def:m1:matching-algorithms-and-implied-spreads:fifo}{price-time priority} and rebates. The same series trades under
both at once, on different exchanges of the same group.

\section{Price-improvement auctions}

\begin{definition}[Price-improvement auction]\label{def:m1:options-market-structure:auction}
In a \emph{price-improvement auction}\index{price-improvement auction} a
member submits a customer's order (the agency order) together with its own
contra order guaranteeing execution of the whole at a stop price at or
better than the best displayed price. The exchange announces the auction,
collects responses for a short fixed period, and executes the agency order
against the best-priced responses first. At the final price the initiating
member retains a guaranteed share.
\end{definition}

\begin{dated}{2026-09}{One exchange's auction, from its rule filings}\label{dat:m1:options-market-structure:aim}
The auction period is set by class, no shorter than 100 milliseconds and no
longer than 3 seconds. If the auction produces no better price, priority
customers resting on the book at the stop price are filled first; the
initiating order then receives up to 50\% of the agency order if one other
participant is at that price, or 40\% if there are two or more; the rest is
shared among the responders. The initiator's share at the final price can
never exceed those percentages of the original order.
\end{dated}

\begin{omfigure}
\begin{tikzpicture}[font=\small]
\draw[thick,-{Stealth[length=2.5mm]}] (0,0) -- (13,0);
\foreach \x/\d/\t in {0.8/$t_0$/{the paired orders\\arrive},4.4/{$t_0$ plus a few ms}/{the auction\\is broadcast},11.2/{$t_0 + 100$ ms}/allocation}{
  \draw[thick] (\x,0.12) -- (\x,-0.12);
  \node[above=3pt,font=\footnotesize\bfseries] at (\x,0.12) {\d};
  \node[below=3pt,font=\footnotesize,align=center] at (\x,-0.12) {\t};}
\foreach \x in {6.2,7.0,7.3,8.6,9.5}{\draw[omThm,thick,-{Stealth[length=1.6mm]}] (\x,0.95) -- (\x,0.08);}
\node[font=\footnotesize,text=omThm] at (7.85,1.2) {responses, each at its own price};
\end{tikzpicture}

\omcaption{A \omterm{def:m1:options-market-structure:auction}{price-improvement auction} at the minimum period. A responder
has a tenth of a second, less its distance from the exchange, to receive the
message, price the option and answer.}\label{fig:m1:options-market-structure:timeline}
\end{omfigure}

The auction gives the customer a price no worse than the stop, sometimes
better; it gives the initiating firm a guaranteed share of flow it selected;
and it gives everyone else a chance to take the order by improving. Whether
it is competitive is an empirical matter that turns on two numbers: how many
firms respond, and how often they improve
(\cref{fig:m1:options-market-structure:auction}).

\begin{omfigure}
\begin{tikzpicture}
\begin{axis}[width=11cm,height=5cm,
  xlabel={probability that a responder improves by one tick},ylabel={initiator's share (\%)},
  tick label style={font=\small},label style={font=\small},
  xmin=0,xmax=1,ymin=0,ymax=75,grid=major,grid style={gray!25},
  legend columns=2,legend style={font=\footnotesize,at={(0.5,-0.3)},anchor=north,draw=none,fill=none,/tikz/every even column/.append style={column sep=8pt}},
  table/col sep=comma]
\addplot[omDef,very thick,mark=*,mark size=1.4pt] table[x=p_improve,y=share_1]{figdata/markets-1/24-options-market-structure/auction.csv};
\addplot[omThm,very thick,dashed,mark=square*,mark size=1.4pt] table[x=p_improve,y=share_3]{figdata/markets-1/24-options-market-structure/auction.csv};
\legend{one responder on average,three responders on average}
\end{axis}
\end{tikzpicture}

\omcaption{Simulated auctions of a 100-lot order: the initiating firm's
average share. With one responder on average who never improves, it keeps two
thirds; with three who improve half the time, an eighth. The guarantee is
worth what the competition lets it be worth. Data: the tutorial's
simulation.}\label{fig:m1:options-market-structure:auction}
\end{omfigure}

\section{Complex orders}

\begin{definition}[Complex order book]\label{def:m1:options-market-structure:complex}
A \emph{complex order book}\index{complex order book} is a book for
multi-leg strategies (spreads, straddles, combinations with the share) quoted
at one net price and executed as a unit. A complex order can trade against
other complex orders or, leg by leg, against the simple books, the exchange
ensuring that every leg prints at a permitted price.
\end{definition}

It is the options counterpart of the \omterm{def:m1:matching-algorithms-and-implied-spreads:calendar}{calendar spread} and its implied
orders (\cref{def:m1:matching-algorithms-and-implied-spreads:implied}): a
vertical spread whose legs are each ten cents wide can be quoted and traded a
cent or two from its fair value as a package, because the package carries a
fraction of the risk of either leg. A large share of institutional options
volume is in such packages, and a \omterm{def:m1:what-a-trading-firm-does:market-maker}{market maker} who quotes only single series
sees only part of the flow.

\section{Increments, fees and the tape}

\begin{definition}[Penny program]\label{def:m1:options-market-structure:penny}
The \emph{penny program}\index{penny program} is the industry-wide rule that
sets the minimum quoting increment of listed options by class: a cent for
the most active classes, five and ten cents for the rest.
\end{definition}

\begin{definition}[Marketing fee]\label{def:m1:options-market-structure:marketing}
A \emph{marketing fee}\index{marketing fee} is a charge levied by an exchange
on \omterm{def:m1:what-a-trading-firm-does:market-maker}{market makers}' trades with customers and pooled to pay brokers for
routing customer orders to that exchange: \omterm{def:m1:retail-flow-and-wholesaling:pfof}{payment for order flow} collected
and distributed through the exchange itself.
\end{definition}

\begin{definition}[OPRA]\label{def:m1:options-market-structure:opra}
\emph{OPRA}\index{OPRA}, the Options Price Reporting Authority, is the
\omterm{def:m1:us-equity-market-structure:sip}{securities information processor} of US listed options: it consolidates the
quotes and trades of all options exchanges into one feed and publishes the
\omterm{def:m1:us-equity-market-structure:nbbo}{national best bid and offer} of every series.
\end{definition}

\begin{dated}{2026-09}{Pennies and messages}\label{dat:m1:options-market-structure:penny}
\emph{Increments.} The penny pilot became the permanent Penny Interval Program
in 2020. All series of three of the most active fund options quote in cents
at any price; other classes in the programme quote in cents below \$3.00 and
in five-cent steps from \$3.00; a class stays in the programme while it is
among the 425 most active. \emph{The feed.} The processor's capacity notice of
September 2025 projects, for July 2026, 311 billion messages a day and a peak
of 13.6 million messages per 100 milliseconds on one stream, 4.4 gigabits in
that tenth of a second, to be doubled for the redundant stream and raised by
a tenth for retransmissions; median latency through the processor is under 18
microseconds.
\end{dated}

\begin{omfigure}
\begin{tikzpicture}
\begin{axis}[width=9.5cm,height=4.4cm,ybar,bar width=14pt,
  ylabel={billion a day},
  symbolic x coords={10/2025,1/2026,7/2026,1/2027,7/2027},xtick=data,
  tick label style={font=\small},label style={font=\small},xticklabel style={font=\small\strut},
  ymin=0,ymax=400,ymajorgrids,grid style={gray!25},
  nodes near coords,nodes near coords style={font=\footnotesize,fill=white,inner sep=1pt},table/col sep=comma]
\addplot[fill=omDef!60,draw=omDef] table[x=date,y=messages_bn_day]{figdata/markets-1/24-options-market-structure/opra.csv};
\end{axis}
\end{tikzpicture}

\omcaption{Projected capacity of the consolidated options feed, in billions
of messages a day, from the processor's notice of September 2025. A firm that
wants every quote of every series must be able to read this. Data: the
notice.}\label{fig:m1:options-market-structure:opra}
\end{omfigure}

A tick of five cents on a three-dollar option is 1.7\% of its price: the
increment is a floor under the spread, a source of \omterm{def:m1:what-a-trading-firm-does:market-maker}{market makers}' revenue,
and the reason \omterm{def:m1:options-market-structure:auction}{price-improvement auctions} have something to improve. The
fee side is as varied as in equities (\cref{ch:m1:exchanges-brokers-venues}),
with one more dimension: fees depend on \emph{who} is trading (customer,
professional customer, firm, \omterm{def:m1:what-a-trading-firm-does:market-maker}{market maker}) as much as on whether the order
makes or takes. A \omterm{def:m1:what-a-trading-firm-does:market-maker}{market maker}'s net price on a trade is known only after the
counterparty's capacity is.

\section{Tutorial: dividing an order}\label{tut:m1:options-market-structure:alloc}

\begin{tutorial}
\textbf{Goal.} Allocate an incoming order under \omterm{def:m1:options-market-structure:customer}{customer priority} with a
designated-maker entitlement; run a \omterm{def:m1:options-market-structure:auction}{price-improvement auction}; measure how
competition erodes the initiator's guarantee. \textbf{End state:} the three
data figures of this chapter.
\begin{tutsteps}
\item \textbf{At the quote.} Customers, then the designated maker, then the
other professionals by size.
\omcode{code/firm/match/firm_optmatch.py}{26}{51}{Customer priority, an entitlement, and pro rata for the rest.}\label{lst:m1:options-market-structure:alloc}
\item \textbf{In the auction.} Better prices first; at the stop, customers,
then the initiator's share, then the responders; the initiator takes what is
left, since it guaranteed the order.
\omcode{code/firm/match/firm_optmatch.py}{96}{120}{The end of a price-improvement auction: what happens at the stop price, after better-priced responses have been filled.}\label{lst:m1:options-market-structure:auction}
\item \textbf{Many auctions.} Poisson responders, each improving with a
given probability.
\end{tutsteps}
\textbf{What to change next.} Let responders bid for part of the order, in
proportion to their confidence, and let the initiator choose a stop price
inside the spread: find the stop that maximises its expected profit when its
edge is six cents at the displayed price.
\end{tutorial}

\section{Build: options allocation}\label{bld:m1:options-market-structure:optmatch}

\begin{build}
\bfield{Purpose} The options venue of the miniature firm's simulated
exchange needs the two mechanisms that distinguish it from the futures
venue of \cref{bld:m1:matching-algorithms-and-implied-spreads:match}.
\bfield{Interface} A resting \texttt{Quote(oid, qty, capacity, designated)};
the function \texttt{customer\_priority\_pro\_rata(book, qty, entitlement\_pct)};
\texttt{price\_improvement\_auction(agency\_qty, stop\_price, side, responses,
book\_customers, one\_other\_pct, many\_others\_pct)} returning the final price,
the fills, the initiator's quantity and whether the price improved.
\bfield{Rules} Integers only; conservation (fills plus the initiator's
quantity equal the order); customers never receive less than under any other
rule; the initiator's share at the final price never exceeds its percentage of
the \emph{original} order when others are present, and is the whole order when
nobody responds.
\bfield{Acceptance tests} \texttt{code/firm/match/tests/test\_firm\_optmatch.py}:
customers first whatever their time; the entitlement against the pro-rata
share; auctions with no, one and several responders; improvement taking the
order away; a selling customer; conservation.
\bfield{Stretch} A \omterm{def:m1:options-market-structure:complex}{complex order book} for two-leg spreads with leg prices
constrained by the simple books.
\end{build}

\begin{omsources}
\item Cboe Exchange, rules of the Automated Improvement Mechanism, as filed
with the SEC (SR-CBOE-2019-045 and later filings, Exhibit 5).
\item US Securities and Exchange Commission, orders on the Penny Interval
Program, 2020 (for example Release 34-89167, Exhibit 5); The Options Industry
Council, ``Why certain options trade in penny increments''.
\item OPRA, \emph{Revised OPRA Capacity Projections}, notice to multicast
data subscribers, 15 September 2025.
\end{omsources}

\section{Exercises}

\begin{exercise}[$\star$]\label{exo:m1:options-market-structure:1}
The level of \cref{fig:m1:options-market-structure:rules} receives an order
for 12 lots under \omterm{def:m1:options-market-structure:customer}{customer priority}. Give the fills.
\end{exercise}

\begin{exercise}[$\star$]\label{exo:m1:options-market-structure:2}
Same level, 115 lots, \omterm{def:m1:options-market-structure:customer}{customer priority} and a 40\% entitlement for the
designated maker. Reproduce the figure's numbers by hand.
\end{exercise}

\begin{exercise}[$\star$]\label{exo:m1:options-market-structure:3}
In a class of the penny programme that is not one of the three all-penny
funds, which of these limit prices are valid: 2.43, 2.99, 3.02, 3.05, 12.10,
12.12? What is the minimum spread, in percent of price, of an option quoted
around 3.10?
\end{exercise}

\begin{exercise}[$\star\star$]\label{exo:m1:options-market-structure:4}
From \cref{dat:m1:options-market-structure:penny}: give the average message
rate per second over a 6.5-hour day, the peak rate per second, and their
ratio. Give the bandwidth a firm must provision to take both streams with the
retransmission margin.
\end{exercise}

\begin{exercise}[$\star\star$]\label{exo:m1:options-market-structure:5}
A 100-lot agency order is auctioned with a stop of 2.50. Give the
initiator's fill when: nobody responds; one firm responds at 2.50 for 100;
three firms respond at 2.50 for 100 each; one firm responds at 2.49 for 60 and
two at 2.50.
\end{exercise}

\begin{exercise}[$\star\star$]\label{exo:m1:options-market-structure:6}
The 50 call is quoted 2.40 at 2.50 and the 55 call 0.90 at 1.00. Give the net
price of buying the 50--55 call spread leg by leg at displayed prices, its
midpoint value, and the width of the package's market. Why can a complex
book quote it much tighter than 0.20?
\end{exercise}

\begin{exercise}[$\star\star\star$]\label{exo:m1:options-market-structure:7}
\emph{Coding.} With \texttt{simulate\_auctions(4000, m, 0.5, 24)} report the
initiator's share and the share of the order that received a better price,
for $m = 1$ and $m = 3$ responders on average. What does the customer gain
from the third responder?
\end{exercise}

\begin{exercise}[$\star\star\star$]\label{exo:m1:options-market-structure:8}
\emph{Find the flaw.} A proprietary firm that is not a registered \omterm{def:m1:what-a-trading-firm-does:market-maker}{market
maker} writes: ``Our 5-lot bid joined the best bid first, ahead of 450 lots of
market-maker quotes. We will be filled first.'' The exchange uses \omterm{def:m1:options-market-structure:customer}{customer
priority} and pro rata. What will happen when 100 lots are sold at that
price?
\end{exercise}

\section{Problem: The Auction}

\begin{problem}[{Weekend problem --- what a guaranteed share is worth}]\label{pb:m1:options-market-structure:1}
A \omterm{def:m1:retail-flow-and-wholesaling:wholesaler}{wholesaler} receives a customer's order to buy 100 calls. The market is 2.40
at 2.50; the \omterm{def:m1:retail-flow-and-wholesaling:wholesaler}{wholesaler} values the call at 2.44. It can start a
\omterm{def:m1:options-market-structure:auction}{price-improvement auction} with a stop of 2.50, under the rules of
\cref{dat:m1:options-market-structure:aim}. From experience: nobody responds
in 10\% of auctions; one firm responds at the stop in 25\%; several respond at
the stop in 35\%; in 20\% responders improve to 2.49 for 60 lots, the rest
trading at the stop with several participants; in 10\% responders improve for
the whole order.

\medskip\noindent\textbf{Part I --- The share.}
\begin{enumerate}
\item Give the \omterm{def:m1:retail-flow-and-wholesaling:wholesaler}{wholesaler}'s fill in each of the five cases.
\item Give its expected fill.
\item Give its edge per lot at 2.50, in dollars per contract.
\item Give its expected gross profit per auction.
\item Give the customer's expected \omterm{def:m1:retail-flow-and-wholesaling:pfof}{price improvement} per contract.
\end{enumerate}

\medskip\noindent\textbf{Part II --- The alternatives.}
\begin{enumerate}[resume]
\item If the \omterm{def:m1:retail-flow-and-wholesaling:wholesaler}{wholesaler} simply routed the order to the exchange's book, what
would the customer pay, and what would the \omterm{def:m1:retail-flow-and-wholesaling:wholesaler}{wholesaler} earn?
\item The \omterm{def:m1:retail-flow-and-wholesaling:wholesaler}{wholesaler} pays the broker 30 cents per contract for the order.
Give its net expected profit per auction.
\item It could instead choose a stop of 2.48. What changes for the customer,
for the \omterm{def:m1:retail-flow-and-wholesaling:wholesaler}{wholesaler}'s edge, and for the responders' incentive to improve?
\item A responder values the call at 2.45. At what prices is responding
worth its while, and what does it need to respond in time?
\end{enumerate}

\medskip\noindent\textbf{Part III --- The \omterm{def:m1:what-a-trading-firm-does:market-maker}{market maker} on the book.}
A \omterm{def:m1:what-a-trading-firm-does:market-maker}{market maker} was offering 50 lots at 2.50 on the book when the auction
began.
\begin{enumerate}[resume]
\item Does its resting quote take part in the auction's allocation at 2.50?
What must it do to take part?
\item Why might it prefer not to respond, and keep its quote for the orders
that are \emph{not} auctioned?
\item Which orders are not auctioned, and what does that imply about the flow
that reaches the displayed quotes?
\item How should that change the width of its displayed quotes?
\item Relate this to \cref{ch:m1:retail-flow-and-wholesaling}.
\end{enumerate}

\medskip\noindent\textbf{Part IV --- Judgement.}
\begin{enumerate}[resume]
\item Is a 100-millisecond auction competitive? Give the argument on each
side.
\item Who pays for the \omterm{def:m1:options-market-structure:marketing}{marketing fee}, in the end?
\item Why do exchange groups operate several options exchanges with different
allocation rules?
\item What would a move to cent increments in every class do to auctions?
\item State the \emph{named result}: the expected allocation to the
initiating firm.
\item In one sentence: what is a customer's option order worth, and to whom?
\end{enumerate}
\end{problem}

\section{Interview questions}

\begin{interviewq}[$\star$ \iqroles{trader, researcher}]\label{iq:m1:options-market-structure:1}
How does the US options market differ in structure from the US equity
market?
\end{interviewq}

\begin{interviewq}[$\star$ \iqroles{trader, developer}]\label{iq:m1:options-market-structure:2}
What is \omterm{def:m1:options-market-structure:customer}{customer priority}, and why do exchanges grant it?
\end{interviewq}

\begin{interviewq}[$\star\star$ \iqroles{trader, researcher}]\label{iq:m1:options-market-structure:3}
Walk me through a \omterm{def:m1:options-market-structure:auction}{price-improvement auction} from the point of view of the
initiator and of a responder.
\end{interviewq}

\begin{interviewq}[$\star\star$ \iqroles{developer}]\label{iq:m1:options-market-structure:4}
You must consume the consolidated options feed. What are the engineering
problems?
\end{interviewq}

\begin{interviewq}[$\star\star$ \iqroles{trader, researcher}]\label{iq:m1:options-market-structure:5}
Why do \omterm{def:m1:options-market-structure:mm}{options market makers} need \omterm{def:m1:options-market-structure:protection}{quote protection}, and what does it cost
the rest of the market?
\end{interviewq}

\begin{interviewq}[$\star\star\star$ \iqroles{researcher, trader}]\label{iq:m1:options-market-structure:6}
You make markets in options on a displayed book. How does the existence of
auctions and \omterm{def:m1:retail-flow-and-wholesaling:wholesaler}{wholesalers} change the flow you see, and what do you do about
it?
\end{interviewq}
